% !TeX TS-program = xelatex
% 虚构演示简历 — 仅用于 hellojobs 测试数据
\documentclass{resume-photo}
\usepackage{xeCJK}
\setCJKmainfont{Noto Serif CJK SC}
\usepackage{fontspec}
\usepackage{enumitem}
\setlist[itemize]{itemsep=0pt, topsep=1pt, parsep=0pt, partopsep=0pt}

\ResumeName{董梅}

\begin{document}

\ResumeContacts{
  138-0020-0020,%
  \ResumeUrl{mailto:dongmei@ruc.edu.cn}{dongmei@ruc.edu.cn},%
  \textnormal{中国人民大学 | 信息管理与信息系统 · 硕士 | 2022—2025}%
}

\ResumeTitle

\noindent 技术型产品经理，熟悉 SQL、埋点体系与 PRD 评审；推动搜索排序策略需求从评审到上线周期缩短 25\%。注重可观测性与文档沉淀，习惯在评审阶段拆解风险；参与线上复盘与跨团队联调，英语 CET-6，可阅读官方文档与 RFC（演示数据）。

\section{教育背景}
\ResumeItem
{中国人民大学}
[\textnormal{信息学院 | 信息管理与信息系统 | 硕士}]
[2022.09—2025.06]

\begin{itemize}
  \item 数字化治理与数据分析方向。
\end{itemize}


\ResumeItem
{科研训练与竞赛}
[\textnormal{大创 / 挑战杯 / 互联网+（演示）}]
[2021—2025]

\begin{itemize}
  \item 主持或核心参与校级大创项目 1 项，负责方案设计、里程碑管理与中期答辩材料。
  \item 挑战杯 / 互联网+ 校赛金奖团队成员，负责技术方案书与 Demo 部署脚本。
  \item 在实验室周会中汇报进展 10+ 次，形成实验记录与可复现实验脚本仓库。
  \item 协助导师撰写专利交底书 1 份（流程中，测试数据）。
\end{itemize}



\section{专业技能}
\begin{itemize}
  \item 熟悉 Axure/Figma、SQL、A/B 实验设计与漏斗分析。
  \item 掌握 Python 数据处理与简单自动化报表。
  \item 了解微服务边界与接口文档（OpenAPI）协作。
  \item 熟悉 Linux 日常运维与 Shell，具备日志检索与 CPU/内存突增初步定位能力。
  \item 掌握 Git / CI 与 Code Review，了解 Docker 与 Kubernetes 基本编排与滚动发布。
  \item 了解 OAuth2 / JWT、接口幂等与 REST 规范，具备单元测试与集成测试习惯（JUnit/pytest 等）。
  \item 能阅读英文技术文档，具备跨团队沟通与需求澄清能力（测试简历）。
\end{itemize}

\section{实习经历}
\ResumeItem
{滴滴出行}
[\textnormal{产品实习生（技术中台方向）}]
[2024.06—2024.12]

\begin{itemize}
  \item \textbf{职责}: 司机端增长活动配置平台需求管理与验收。
  \item \textbf{成果}: 需求返工率从 18\% 降至 9\%。
\end{itemize}


\ResumeItem
{美团 | 到店技术部}
[\textnormal{后端开发日常实习}]
[2023.07—2024.01]

\begin{itemize}
  \item \textbf{背景}: 营销活动报名接口在晚高峰出现线程池打满与超时尖刺。
  \item \textbf{方案}: 引入异步化改造与 Sentinel 限流，拆分热点读请求至本地缓存。
  \item \textbf{指标}: 接口 P99 从 890ms 降至 260ms，错误率从 0.8\% 降至 0.05\%。
  \item \textbf{工程}: 补齐监控大盘与告警分级，编写上线 checklist 与回滚脚本。
  \item \textbf{协作}: 与前端、测试对齐幂等键与重试策略，完成灰度 5\%→100\%。
\end{itemize}



\section{项目经历}
\ResumeItem
{指标体系白皮书整理}
[\textnormal{校内课题}]
[2024.02—2024.06]

\begin{itemize}
  \item 梳理 120+ 核心指标口径，被课题组采纳为基线文档。
\end{itemize}


\ResumeItem
{日志检索与告警平台}
[\textnormal{ELK + 规则引擎（演示）}]
[2023.10—2024.03]

\begin{itemize}
  \item \textbf{痛点}: 多业务线日志格式不统一，检索慢且告警噪声高。
  \item \textbf{落地}: 制定 log schema，Filebeat 多管道 + Ingest Pipeline 清洗。
  \item \textbf{指标}: 检索 P95 从 4.2s 降至 0.9s，告警误报率下降 55\%。
  \item \textbf{运营}: 建立 on-call 轮值与告警分级 SOP。
\end{itemize}



\section{个人荣誉}
\begin{itemize}
  \item 中国人民大学 优秀学生干部
  \item 互联网+ 校赛 金奖
  \item 全国计算机等级考试四级（网络工程师）（演示）
  \item 校级「优秀学生干部 / 优秀团员」或技术社团骨干（综合测评前 15\%）
  \item 开源贡献：向知名项目提交被合并的 PR（文档/测试/feature）（演示）
\end{itemize}

\end{document}
